% A.70 candidate: notes only, unregistered during checkpoint freeze.
\subsection{Every two-power dihedral geometric closure}

\begin{statement}{A.70}{Proved}{Both cusp assertions for every two-power dihedral closure}
Let \(Y\to U=E\setminus\{O\}\) be an actual connected finite
\'{e}tale cover whose geometric Galois closure has dihedral group
\(D_{2m}\), where \(m\) is a power of two. Then both assertions
of Task A.5 hold for \(Y\), at every coefficient prime, every
induced cusp packet, and every permitted open derived subgroup.
In particular its full fixed geometric pro-\(\ell\) subgroup is
its selected cusp inertia, and every nontrivial arithmetically
normalized procyclic pro-\(\ell\) subgroup is cuspidal.

For \(m\ge4\), there are exactly three geometric Galois covers
with this group. As actual covers of the fixed elliptic base they
are the three models
\[
 C_{m,a,b}:\ v^2=w^{2m}+aw^m+b,\qquad
 (a,b)\in\{(0,-1),(6,1),(-6,1)\},                 \tag{A.70.1}
\]
with the maps below. Every intermediate of each is included.
The order-four case has one geometric \(V_4\) cover and is
handled separately. This conclusion does not cover arbitrary
nilpotent closures or products of two nonabelian Sylow groups;
the universal statement A.45 remains open.
\end{statement}

\begin{proof}
\paragraph{Three actual rotation quotients.}
For each model in (A.70.1), put
\[
 t=w^m,\qquad u=w^{m/2}v,\qquad
 E_{a,b}:u^2=t^3+at^2+bt.                         \tag{A.70.2}
\]
This is a degree-\(m\) map. Indeed the function-field extension
over \(\overline K(t)\) through \(\overline K(w,v)\) has
degree \(2m\), while \(\overline K(E_{a,b})\) has degree two.
The polynomial in (A.70.1) has simple roots, and is not a square
in \(\overline K(w)\). Thus the compact source is connected,
smooth after normalization, and has genus \(m-1\).
The cyclic rotation group is generated by
\[
                 r(w,v)=(\zeta_mw,-v).            \tag{A.70.3}
\]
It fixes \((t,u)\) and has order \(m\), so (A.70.2) is its
actual quotient, not merely an isogenous elliptic factor.

The degree-two isogeny with kernel \(\{O,(0,0)\}\) on
\(E_{a,b}\) is given by
\[
 t'=t+a+b/t,\qquad u'=u(1-b/t^2),\qquad
 u'^2=t'^3-2at'^2+(a^2-4b)t'.                    \tag{A.70.4}
\]
For \(a=\pm6,b=1\), identify its target with the fixed
\(E:Y^2=X^3-X\) by
\[
                  X=(t'-2a/3)/4,\qquad Y=u'/8.
                                                               \tag{A.70.5}
\]
For \((a,b)=(0,-1)\), use
\[
                  X=-it'/2,\qquad Y=-(1+i)u'/4.
                                                               \tag{A.70.6}
\]
These are the explicit isogenies already checked in A.4. Their
composites with (A.70.2) have degree \(2m\) and have exactly
four points above \(O\), namely the two points above each of
\(w=0,\infty\). The cyclic map is ramified only there, with
index \(m/2\), because \(\operatorname{div}(t)=2(0,0)-2O\)
on \(E_{a,b}\). The isogeny is unramified.

Choose \(\eta\) with \(\eta^m=b\), after a finite extension.
The actual reflection
\[
 s(w,v)=\left(\eta/w,\ -\frac{b}{\eta^{m/2}}\frac v{w^m}\right)
                                                               \tag{A.70.7}
\]
satisfies \(s^2=1\) and \(srs=r^{-1}\). It induces
\[
                   (t,u)\longmapsto(b/t,-bu/t^2),
                                                               \tag{A.70.8}
\]
which is translation by \((0,0)\) on \(E_{a,b}\).
Thus the composite to the fixed base is invariant under a faithful
\(D_{2m}\) action of order equal to its degree, and is the actual
Galois cover claimed in the statement.

\paragraph{The three degree-two covers of the fixed base are distinct.}
The nonzero kernel point of the dual of (A.70.4) is \((t',u')=(0,0)\).
One can see this directly: the two nonkernel nonzero two-torsion
points on \(E_{a,b}\) satisfy \(t^2+at+b=0\), so both map
to \((0,0)\); the pullback--dual identity says that this image
is precisely the dual kernel. This is the usual actual dual-isogeny
identity \cite[III.6.1--6.2]{fghSilvermanAEC}.
Under (A.70.5)--(A.70.6) these
three dual kernels on the \emph{fixed} elliptic curve are
\[
 \begin{array}{c|c}
 (a,b)&\text{nonzero dual-kernel point on }E\\ \hline
 (0,-1)&(0,0)\\
 (6,1)&(-1,0)\\
 (-6,1)&(1,0).
 \end{array}                                                   \tag{A.70.9}
\]
An isomorphism of degree-two covers over the identity of \(E\)
preserves this dual kernel. Hence the three rotation quotients
are different covers of the specified base. The two nonstable
source curves may be isomorphic over an algebraic extension;
that isomorphism covers \([i]\) on \(E\), not its identity,
and does not identify their marked covering kernels.

\paragraph{Exhaustion of the geometric dihedral closures.}
The geometric fundamental group of \(U\) is free profinite on
two generators. Connected geometric Galois covers with group
\(D_{2m}\) correspond to continuous epimorphisms onto this
finite group, modulo its automorphism group. Equivalently they
are generating ordered pairs modulo automorphisms.
For \(m\ge4\), the rotation subgroup \(R=\langle r\rangle\)
is characteristic: it is generated by the elements of order
strictly greater than two. Every automorphism has the form
\[
       r\longmapsto r^e,\qquad s\longmapsto r^f s,
       \qquad e\in(\mathbf Z/m)^\times,\ f\in\mathbf Z/m.
                                                               \tag{A.70.10}
\]
A generating pair must have at least one reflection. If only
one entry is a reflection, the other is an odd rotation. If
both entries are reflections, their exponents differ by an
odd integer. Formula (A.70.10) is transitive on each of the
three resulting types, with representatives
\[
                      (r,s),\qquad(s,r),\qquad(s,rs).
                                                               \tag{A.70.11}
\]
These types cannot be equivalent because their images in the
characteristic quotient \(D_{2m}/R\simeq C_2\) are the three
different nonzero ordered binary pairs. Thus there are exactly
three geometric covering kernels, distinguished by the actual
rotation quotient. The three models above have distinct such
quotients, so they exhaust these kernels.

For \(m=2\), the group is \(V_4\), its rotation subgroup is
not characteristic, and \(\operatorname{Aut}(V_4)\) is transitive
on generating ordered pairs. There is one cover, the maximal
exponent-two abelian cover, which extends to the elliptic
isogeny \([2]\). Its intermediates have genus one and are
settled by A.40. A cyclic group of order two, if included under
the degenerate dihedral notation, is likewise an elliptic-isogeny
case.

\paragraph{Every intermediate, including the stable family.}
For the two nonstable models, A.69 already proves both assertions
for every actual intermediate. Its potential-good and cusp-torsion
statements are intrinsic to the marked compact curve, so they
apply to the explicit maps (A.70.5), including the plus-six map.
For the stable model, A.44 proves potential good reduction of its
Jacobian at two, special \(2\)-rank zero, and annihilation of
every cusp difference by \(4m\). We spell out why the same
hypotheses reach every intermediate, rather than infer fixedness
through a pro-prime quotient.

At every cusp of (A.70.1), the actual inertia is
\(I=\langle r^2\rangle\), normal in \(D_{2m}\).
For every proper subgroup \(H<D_{2m}\), the quotient curve
\(D=C_{m,0,-1}/H\) has cusp count
\[
                   [D_{2m}:HI]\ge2.              \tag{A.70.12}
\]
Indeed \(I=\Phi(D_{2m})\); more elementarily, \(HI=D_{2m}\)
would supply an odd rotation and a reflection in \(H\), forcing
\(H=D_{2m}\). The actual norm map makes \(J_D\) an isogeny
quotient of the source Jacobian, so it has potential good
reduction with special \(2\)-rank zero. Over an algebraically
closed field the pushforward of a point is its image with
coefficient one. Norms of the actual functions annihilating
source cusp differences therefore kill every difference of
target cusps by \(4m\). The compact genus of \(D\) is at least
one, because it has a nonconstant finite map to \(E\).

Consequently the torsion-divisor criterion in A.41 gives a
rational joining path to a distinct companion for every cusp
of \(D\). A.42 proves full fixed inertia at every coefficient
prime; A.59 proves every normalizer assertion, since the actual
closure remains a two-group. The quotient \(H=D_{2m}\) is the
original \(U\), already settled.

Finally any cover in the statement becomes, after a finite
number-field extension, an actual intermediate of one of these
three geometric closures. The covering isomorphisms, subgroup
and marked points can all be defined over a finite extension.
A geometric path used to identify base points need not be
defined over that field: it gives the arithmetic comparison up
to geometric inner conjugacy, under which the packet, fixed-group
and normalizer assertions are invariant.
The finite-field-restriction argument of A.41--A.42
then gives the original permitted open derived subgroups and
packets; normalizer recognition also descends by intersecting
its open arithmetic image with this finite-index Galois subgroup.
No assertion that fixed points commute with a quotient is used.
The exact isogeny identities and finite generating-pair support
are retained in \path{notes/fgh_A70_certificate.py}.
\end{proof}
